\section{The denominator bound on the selected design}

% BEGIN G87D1
\paragraph{G87D1. Exhaustion of the good eigenspaces on the selected members.}
Work on the eventual layer range of F86D2, with its original
coefficient embedding, level tuple and norm. For an arbitrary good
point $x$ and nonzero upper vector $f\in E_x$, write $y=Af$ in
the linked realization. M87U1 gives its unramified eigenvalues
with their distinguished string and M87U2 gives its good projector.
D87C1--3 furnish one scalar component and one polynomial to which
the doubling calculation applies; D87C4 identifies all the other
data, including the finite excluded set and the quotient measure.
The compact archimedean factor and the excluded-place factor are
nonzero. By Y87R3 the doubling integral has a simple pole with
nonzero residue. If all the test functions gave zero downward
lift of this scalar component, Y87R4--5 would force that residue
to vanish. Thus the original $y$ has a nonzero varying-test
integral.

Apply Q87F7 to this same vector: its prescribed downward theta
lift is nonzero. This works for every nonzero $f$ in the eigenspace,
including arbitrary linear combinations across its multiplicities.
M87U3 gives $T_iA_x=E_x$ for every $i,x$, with $A_x$ the lower
parameter space. Thus both remaining original fields
$\mathrm{what}_2$ and $\mathrm z$ hold, on the same members
as the other sixteen fields and the four compatibility ties.
This gives the eventual layer assertion for their conjunction.
The indices $x$ here label Hecke points on a member. No
archimedean vector, local Schwartz test or constant from the
pole calculation must be chosen uniformly across members:
these arguments prove an exact kernel statement at each of them.
All level and character parameters in constructing the member
retain the uniform quantifiers of F86D2.
% END G87D1

% BEGIN G87D2
\paragraph{G87D2. The original global denominator estimate.}
Take the built class $\mathcal A$ of the common construction and
put $\mathcal C=\mathrm{WithChN}(\mathcal A)$. In the selected
proof graph the denominator assertion is
$\mathrm{Closed}(\mathrm{Des}(\mathcal C,3),\mathrm{LinkN})$.
By G87D1 every original input of the member implication in
A85D3 holds on the layer range of X87S4. Apply X87S3 before the
single joint class selection. X87S1--2 and X87S4 then give the
member assertion at the norm of each actual selected object, for
all its guarded indices. This installs the denominator row in
that joint construction without assuming the row as an input.
The separate existence question for the common design supplies
its families to all the rows at once.

Take $\varepsilon_1=1$. For each $0<\varepsilon\le1$ choose the
transport parameter $K=(0,0,0,0)$ as in X86B2, before the family
and height constant. Fix $\mathcal F$ in the stated class and
$C_H$ satisfying its original $\mathcal F.\mathrm{Large}(K,C_H)$
condition. Membership gives the twist condition. On this
family take $C_T=40$ as in A85D3. The original tame lifting estimate
at that constant, the cut comparison P84B2--3 and the common
family bound of A83G2--3 give $C,x_0$ before any late scale,
index or lift is chosen. One may enlarge $x_0$ to exceed
$\mathcal F.x_1$ as in X86B1. For every $x\ge x_0$ and family index
$i$, put $D=\mathcal F.D(x,i)$ and $L=\mathcal F.L(x,i)$.
Under $\mathrm{ShaFinOdd}(D)$, for every permitted $j$, summability
of $v_\ell=\mathrm{vvTerm}(D,\ell)$ gives
\[
 \mathrm{etaSum}(L,j)\ \le\ dZ_0(j)\!
 \left[C\left(fm\,b L_{\rm sc}^{2}
                +(\log x)^2 N_c\right)+\sum_\ell v_\ell\right].
\]
Here $f,m$ are the member's rank parameters, $b,L_{\rm sc},N_c$
come from its size datum and $dZ_0(j)$ is its original lift rank.
The sum on the left has weight $\mathrm{lwP}(w)e(w)$ and value
$\eta_J(j,w)$ at each supported $w$ with $\mu_{\rm At}(w)\ne\mu_0$.
On the right $v_\ell$ is zero at nonprimes and is
$\log\ell$ times the sum of the indicated cohomological lengths
at $1\le e<t$ for the same $D$ at primes $\ell$. The finiteness
guard retains $\ell=2$. The rank multiplier applies to the
whole bracket. The transfer off the cut pays the error
$dZ_0(j)\log N_G$ within the original polynomial allowance as
in P84B2; it does not change the cohomological series.
These are all the clauses of $\mathrm{LinkN}$ with its original
constants and guards. Restriction to the intersection with the
target class is the projection in X87S4. This closes the shared
denominator obligation in the reconstruction for RH and for
unrestricted Dirichlet GRH.
% END G87D2
